\theory{Number theory}{How can we understand the exact patterns hidden in whole numbers, especially in prime numbers?}{Number theory studies the integers, divisibility, remainders, primes, equations, and the ways numbers can be approximated or represented. It begins with familiar arithmetic but reaches algebra, geometry, probability, and computing. Its central method is to replace a large numerical question by a structural one: instead of testing every possibility, factor a number, work with its remainder, or study the shape of all solutions.}{Public-key cryptography, error-correcting codes, fast computer arithmetic, calendars, random-number generation, digital signatures, and the analysis of patterns in data.}{Divisibility, congruences, primes, and Diophantine equations organize the arithmetic structure of integers and related rings.}

\paragraph{The question and its history}
The objects of number theory are numbers such as $0,1,2,\ldots$, the negative integers, fractions, and numbers obtained from polynomial equations. A typical question is whether an equation such as $x^2+y^2=z^2$ has solutions, or whether a very large integer is prime. These questions are exact: a numerical experiment can suggest a pattern, but a proof must explain why the pattern continues forever \cite{number_theory_ref1}.

The subject has several connected parts. Elementary number theory studies divisibility, primes, congruences, Diophantine equations, continued fractions, and partitions. Analytic number theory uses limits and complex functions to study how primes are distributed. Algebraic number theory enlarges the integers to number systems in which polynomial equations become easier to understand. Geometric and arithmetic geometry study integer and rational solutions using curves and higher-dimensional shapes. Probabilistic, combinatorial, computational, and applied number theory use randomness, counting, algorithms, and applications \cite{number_theory_ref1,number_theory_ref2}.

Babylonian tablets record calculations involving right triangles. Greek mathematicians made the theory deductive: Euclid proved that every integer has one prime factorization and that there are infinitely many primes. Diophantus studied equations in integers and rational numbers. The Chinese remainder theorem was developed in ancient China, while Indian mathematicians such as Aryabhata, Brahmagupta, and Bhaskara developed methods for indeterminate equations. In the seventeenth and eighteenth centuries Fermat, Euler, Lagrange, and Legendre created much of modern elementary number theory. Gauss organized congruences and quadratic forms, and nineteenth-century mathematicians connected numbers with algebra and complex analysis \cite{number_theory_ref1,number_theory_ref3}.

\end{historylist}
\section*{3. Theory}
\subsection*{Core theory}
\paragraph{The basic tools}
If $a$ and $b$ are integers, we say that $a$ divides $b$ when $b=ak$ for another integer $k$. Division with remainder says that for $a>0$ there are unique integers $q$ and $r$ with
\[
b=aq+r,\qquad 0\leq r<a.
\]
Repeated division gives the Euclidean algorithm for the greatest common divisor. For example,
\[
252=2(105)+42,\quad 105=2(42)+21,
\]
so $\gcd(252,105)=21$. Reversing these steps also writes the greatest common divisor as an integer combination of the original numbers. This is the extended Euclidean algorithm, and it produces modular inverses \cite{number_theory_ref1}.

The fundamental theorem of arithmetic says that every integer greater than $1$ can be written as a product of primes, and that the product is unique apart from the order of the factors. Thus $84=2^2\cdot3\cdot7$ in only one essential way. This theorem turns divisibility into bookkeeping with prime exponents. It also explains why the greatest common divisor takes the smaller exponent of each prime and the least common multiple takes the larger exponent \cite{number_theory_ref1}.

Congruence notation records remainders. We write $a\equiv b\pmod m$ when $m$ divides $a-b$. The statement $3x\equiv1\pmod7$ has solution $x\equiv5\pmod7$, because $3\cdot5=15\equiv1\pmod7$. A number has a multiplicative inverse modulo $m$ precisely when it is relatively prime to $m$. Congruences allow a difficult calculation to be reduced to a small finite table while preserving the information needed for divisibility \cite{number_theory_ref1}.

\paragraph{Equations and patterns}
A Diophantine equation asks for integer or rational solutions. The equation $ax+by=c$ has an integer solution exactly when $\gcd(a,b)$ divides $c$. For example, $6x+15y=9$ is solvable because $\gcd(6,15)=3$ divides $9$. The solutions form an infinite arithmetic family once one solution is known. Equations such as $x^2+y^2=z^2$, Pell's equation $x^2-Dy^2=1$, and Fermat-type equations require more specialized ideas \cite{number_theory_ref1,number_theory_ref3}.

Continued fractions approximate a real number by fractions while controlling the error. The expansion of $\sqrt2$ begins
\[
\sqrt2=[1;2,2,2,\ldots],
\]
whose convergents include $1$, $3/2$, $7/5$, and $17/12$. These are unusually good approximations. Continued fractions therefore solve approximation problems and also provide a way to analyze quadratic irrational numbers and Pell equations \cite{number_theory_ref1}.

Prime numbers are the indivisible building blocks of the integers. Their abundance is described asymptotically by the prime number theorem:
\[
\pi(x)\sim\frac{x}{\log x},
\]
where $\pi(x)$ counts primes at most $x$. The formula does not identify the next prime, but it predicts the density of primes for large $x$. Riemann's zeta function and Dirichlet $L$-functions give analytic tools for studying finer questions, including primes in arithmetic progressions. The Riemann hypothesis is a famous unproved conjecture about the zeros of the zeta function \cite{number_theory_ref2}.

\paragraph{How the theory solves practical problems}
Modern cryptography uses number theory because multiplication of large primes is fast, while recovering the primes from their product can be difficult. In RSA, a public key contains a large composite number $n=pq$ and an exponent; a private key is built from the hidden factors $p$ and $q$. Modular exponentiation encrypts and decrypts messages, while Euler's theorem explains why the operations undo one another for permitted messages \cite{number_theory_ref2}.

The Chinese remainder theorem combines separate remainder conditions. If $x\equiv2\pmod3$ and $x\equiv3\pmod5$, then $x\equiv8\pmod{15}$. This is useful in computer arithmetic because a large calculation can be performed modulo several smaller, coprime numbers and then recombined. Error-correcting codes use finite fields and polynomial arithmetic to add redundancy, allowing a receiver to reconstruct data after some symbols are changed \cite{number_theory_ref2}.

Elliptic curves provide another example. An equation such as $y^2=x^3+ax+b$ defines a curve with a special addition rule for its points. The arithmetic of these points supports digital signatures and key exchange. The security is not based on multiplication alone but on the difficulty of reversing a point-multiplication operation on a suitably chosen curve \cite{number_theory_ref3}.

\section*{4. Important theorems and results}
\begin{enumerate}[leftmargin=*,itemsep=0.35em]
\item \textbf{Euclid's theorem on primes.} There are infinitely many primes. If $p_1,\ldots,p_k$ were all the primes, the number $p_1p_2\cdots p_k+1$ would have a prime factor not on the list.

\item \textbf{Fundamental theorem of arithmetic.} Every integer greater than $1$ has one prime factorization.

\item \textbf{Bézout's identity.} There exist integers $u,v$ such that $ua+vb=\gcd(a,b)$. This is the reason the extended Euclidean algorithm finds inverses.

\item \textbf{Chinese remainder theorem.} Compatible congruences with pairwise coprime moduli have one solution modulo the product of the moduli.

\item \textbf{Fermat's little theorem and Euler's theorem.} If $p$ is prime and $p$ does not divide $a$, then $a^{p-1}\equiv1\pmod p$. More generally, $a^{\varphi(n)}\equiv1\pmod n$ when $\gcd(a,n)=1$.

\item \textbf{Prime number theorem.} The proportion of integers near $x$ that are prime is approximately $1/\log x$. It describes average distribution, not an exact list.
\end{enumerate}
\noindent\emph{Source for these results:} \cite{number_theory_ref1}.

\theoryapplications
\theoryconnections{number_theory}{%
\item Integer arithmetic supplies divisibility and congruence, while algebra turns those rules into groups, rings, and fields that can be manipulated systematically.
\item Proof identifies which patterns are universal, and analysis becomes necessary when infinitely many integers are counted or primes are studied through generating functions.
\item For example, reducing modulo $m$ turns an infinite search for solutions into a finite search among residue classes, while a zeta function packages information about all primes into one analytic object \cite{number_theory_ref1,number_theory_ref2}.
}{%
\item Number theory helps coding theory by providing finite-field arithmetic in which errors can be detected and corrected.
\item Cryptography uses modular exponentiation together with problems such as factoring or discrete logarithms that are easy to perform in one direction but difficult to reverse.
\item Algebraic geometry and probability borrow arithmetic counts and congruence information, while mathematical physics uses lattices, modular forms, and spectral sums to organize discrete states \cite{number_theory_ref1,number_theory_ref2}.
}
\section*{7. Sample Questions}
\emph{Exercise reference:} \cite{number_theory_ref1}. The cited edition is the source for the exercise set; consult its Exercises section for the official statement and numbering.
\begin{enumerate}[leftmargin=*,itemsep=0.9em]
\item \textbf{Exercise 1. Compute $\gcd(252,105)$ and express it as a combination.} The Euclidean algorithm gives $252=2(105)+42$ and $105=2(42)+21$. Hence the gcd is $21$. Since $21=105-2(42)=105-2(252-2(105))=5(105)-2(252)$, one combination is $21=-2(252)+5(105)$.

\item \textbf{Exercise 2. Solve $14x\equiv8\pmod{22}$.} The gcd of $14$ and $22$ is $2$, which divides $8$. Divide by $2$ to get $7x\equiv4\pmod{11}$. Since $7^{-1}\equiv8\pmod{11}$, $x\equiv32\equiv10\pmod{11}$. The solutions modulo $22$ are $x\equiv10,21\pmod{22}$.

\item \textbf{Exercise 3. Solve $x\equiv2\pmod3$ and $x\equiv3\pmod5$.} Write $x=2+3k$. Then $2+3k\equiv3\pmod5$, so $3k\equiv1\pmod5$. The inverse of $3$ modulo $5$ is $2$, giving $k\equiv2\pmod5$ and $x\equiv8\pmod{15}$.

\item \textbf{Exercise 4. Prove that $\sqrt2$ is irrational.} Suppose $\sqrt2=a/b$ in lowest terms. Then $a^2=2b^2$, so $a$ is even, say $a=2c$. Substitution gives $b^2=2c^2$, so $b$ is even too, contradicting that $a/b$ was reduced. Thus no fraction equals $\sqrt2$.

\item \textbf{Exercise 5. Find the positive integer solutions of $3x+5y=1$.} Since $3(2)+5(-1)=1$, one solution is $(2,-1)$. The general solution is $x=2+5t$ and $y=-1-3t$ for integers $t$, because changing $x$ by $5$ and $y$ by $-3$ preserves the left side.

\end{enumerate}
\section*{8. References}
\begin{thebibliography}{9}
\bibitem{number_theory_ref1} G. H. Hardy and E. M. Wright, \emph{An Introduction to the Theory of Numbers}.
\bibitem{number_theory_ref2} J. H. Silverman, \emph{A Friendly Introduction to Number Theory}.
\bibitem{number_theory_ref3} J. Neukirch, \emph{Algebraic Number Theory}.
\end{thebibliography}
